%======================================================================
% GESTIÓN DE USUARIOS
%======================================================================

\chapter{Gestión de Usuarios}

La seguridad informática es una disciplina que se encarga de proteger la integridad y la privacidad de la información almacenada en un sistema informativo. Un sistema seguro debe ser íntegro, confidencial, irrefutable y tener buena disponibilidad.

La seguridad por diseño (Security by Design) es un paradigma que integra la seguridad en todas las fases del desarrollo del software, desde la concepción inicial hasta el despliegue y mantenimiento.

\section{Crear Usuarios}

La gestión de usuarios permite:

\begin{itemize}
    \item Crear nuevos usuarios con diferentes niveles de acceso
    \item Asignar roles y permisos específicos
    \item Configurar autenticación multifactor
    \item Establecer políticas de contraseña
\end{itemize}

\section{Roles y Permisos}

El sistema implementa un modelo de control de acceso basado en roles (RBAC):

\begin{itemize}
    \item \textbf{Administrador}: Acceso completo a todas las funciones
    \item \textbf{Operador}: Monitoreo y operaciones de lectura/escritura
    \item \textbf{Auditor}: Acceso de solo lectura a logs y reportes
    \item \textbf{Usuario}: Acceso básico según configuración
\end{itemize}

\section{Autenticación}

\subsection{Seguridad Basada en Tokens (JWT)}

El sistema implementa un robusto mecanismo de autenticación basado en JSON Web Tokens:

Un JSON Web Token (JWT) es una cadena codificada en Base64 formada por tres partes separadas por un punto: Header (algoritmo y tipo), Payload (datos del usuario, expiración, roles), y Signature (firma digital para validación).

El estándar JWT fue definido en RFC 7519 en 2015, proporcionando un método stateless para transmitir claims entre partes de forma segura.

Componentes del Token JWT:

\begin{enumerate}
\item \textbf{Header}: Especifica el algoritmo de firma (HS256/RSA) y el tipo de token.
\item \textbf{Payload}: Contiene los claims del usuario (sub, roles, permisos, expiración).
\item \textbf{Signature}: Verifica que el token no haya sido modificado y fue emitido por el servidor.
\end{enumerate}

El flujo de autenticación funciona de la siguiente manera:

\begin{enumerate}
\item El usuario envía credenciales (usuario/contraseña) al servidor.
\item El servidor valida las credenciales y genera un JWT con la información del usuario.
\item El cliente almacena el JWT (localStorage o cookies seguras).
\item En cada request posterior, el cliente envía el JWT en el header Authorization.
\item El servidor valida la firma y los claims antes de procesar la solicitud.
\end{enumerate}

Para mayor seguridad, se recomienda validar también la IP del cliente al generar el token, detectando así intentos de robo de sesión.

\subsection{Autenticación Multifactor (MFA)}

El sistema implementa MFA mediante TOTP:

El TOTP (Time-based One-Time Password) es un algoritmo que genera contraseñas de un solo uso basadas en el tiempo, típicamente intervalos de 30 segundos, utilizado ampliamente en la autenticación multifactor.

La autenticación multifactor incrementa significativamente la seguridad, ya que requiere múltiples formas de verificación independientes.

Los usuarios deben validar su identidad no solo con credenciales tradicionales, sino también mediante un código dinámico de 6 dígitos generado por aplicaciones de autenticación estándar (como Google Authenticator), fortaleciendo la barrera contra el acceso no autorizado y el robo de identidad.

\clearpage
